\section{K2 es el Chogori: modelo base, eficiencia y capacidad Agentic}

Esta sección entra en K2. Yang Zhilin (杨植麟) divide el objetivo de K2 en varias partes: primero, que sea un modelo base muy bueno; segundo, aprovechar al máximo cada porción de datos, es decir, aumentar la token efficiency; tercero, que tenga buena capacidad Agentic y pueda pasar de ser un «cerebro en una cubeta» al uso de herramientas en múltiples turnos; cuarto, que el código abierto lleve esas capacidades al ecosistema.

\begin{figure}[H]
\centering
\includegraphics[width=0.96\textwidth]{k2-design-goals.png}
\caption{Objetivos de diseño de Kimi K2: modelo base, token efficiency, capacidad Agentic y código abierto deben cumplirse a la vez. Diagrama conceptual propio, elaborado a partir de la conversación entre 00:24:58 y 00:54:08.}
\end{figure}

\begin{knowledgebox}{Lectura de la figura: K2 no es una optimización aislada}
El modelo base determina los cimientos; la token efficiency determina cuánta capacidad puede surgir de los mismos datos; Muon y Rephrase son medios para mejorar la eficiencia del entrenamiento; la capacidad Agentic determina si el modelo puede usar herramientas, y el código abierto determina la difusión en el ecosistema.
\end{knowledgebox}

\subsection{Token efficiency: con los mismos datos, el cerebro crece más}

Después del muro de datos, la token efficiency se vuelve más importante. En palabras de Yang Zhilin, con la misma cantidad de datos se busca que el «cerebro» crezca más. K2 aplica Rephrase a los datos y también presta atención a líneas de trabajo como el optimizador Muon. Estas técnicas no buscan solo mejorar en los benchmarks, sino aumentar el rendimiento de entrenamiento de cada token a medida que los datos de alta calidad se vuelven más escasos.

\begin{figure}[H]
\centering
\includegraphics[width=0.96\textwidth]{token-efficiency-loop.png}
\caption{Ciclo de la Token Efficiency: la reescritura de datos, el optimizador y la estrategia de entrenamiento aumentan el rendimiento de cada token. Diagrama conceptual propio, elaborado a partir de la conversación entre 00:27:28 y 00:32:44.}
\end{figure}

\begin{knowledgebox}{Términos clave: conceptos relacionados con la eficiencia del entrenamiento}
\begin{tabularx}{\textwidth}{p{0.22\textwidth}p{0.32\textwidth}X}
\toprule
Concepto & Función & Riesgo \\
\midrule
Token efficiency & Obtener más capacidad con la misma cantidad de tokens & Es difícil de medir de forma aislada y suele confundirse con otras variables. \\
Rephrase & Reescribir los datos para aumentar la variedad de expresiones & Una reescritura de mala calidad introduce ruido. \\
Muon & Línea de trabajo en optimizadores que puede mejorar la eficiencia del entrenamiento & El entrenamiento puede volverse inestable; en la entrevista se menciona que «explota». \\
Data mixture & Proporción de los datos y currículo de entrenamiento & Una proporción equivocada puede dañar la capacidad o la generalización. \\
\bottomrule
\end{tabularx}
\end{knowledgebox}

\subsection{El reto de generalización de los modelos Agentic}

La sección anterior trató la eficiencia del entrenamiento de K2; esta pasa al problema central de la capacidad Agentic: la generalización. Para un modelo Agentic, el mayor reto es generalizar. Si las herramientas, las tareas y la retroalimentación del entorno de entrenamiento son demasiado limitadas, el modelo puede desempeñarse bien solo en escenarios conocidos; lo verdaderamente valioso es que, ante herramientas, tareas y entornos nuevos, siga siendo capaz de explorar en múltiples turnos, invocar herramientas y corregir su comportamiento.

\begin{figure}[H]
\centering
\includegraphics[width=0.94\textwidth]{agentic-generalization.png}
\caption{El reto de generalización de los modelos Agentic: saber usar herramientas no equivale a generalizar a tareas nuevas. Diagrama conceptual propio, elaborado a partir de la conversación entre 00:32:50 y 00:38:00.}
\end{figure}

\begin{importantbox}{Cómo juzgar la generalización Agentic}
Para saber si un modelo es verdaderamente Agentic no basta con ver si invoca una herramienta fija; hay que ver si, ante tareas desconocidas, puede entender la documentación de las herramientas, planificar acciones de varios pasos, corregirse a partir de la retroalimentación y transferir su capacidad a entornos nuevos.
\end{importantbox}

\begin{knowledgebox}{Tres niveles de prueba de la generalización Agentic}
\begin{tabularx}{\textwidth}{p{0.22\textwidth}p{0.34\textwidth}X}
\toprule
Nivel & Qué se evalúa & Modo de falla \\
\midrule
Generalización de herramientas & Si entiende API, comandos y documentación nuevos & Solo sabe usar las herramientas vistas en el entrenamiento. \\
Generalización de tareas & Si puede descomponer objetivos nuevos y dependencias de largo alcance & El plan se ve bien, pero no se puede ejecutar. \\
Generalización de la retroalimentación & Si puede corregirse después de recibir un error & Persiste en el error o repite la misma acción. \\
\bottomrule
\end{tabularx}
\end{knowledgebox}

\subsection{Resumen de la sección}

El núcleo de K2 es equilibrar varios objetivos: un modelo base sólido, eficiencia de entrenamiento, generalización Agentic y un ecosistema de código abierto deben cumplirse a la vez. Una debilidad en cualquiera de estos eslabones limita su paso de modelo a sistema.
